% Generated draft: requires mathematical and source review.
\documentclass[11pt]{article}
\usepackage[margin=25mm]{geometry}
\usepackage{fontspec}
\setmainfont{Latin Modern Roman}
\usepackage{amsmath,amssymb,mathtools,unicode-math}
\setmathfont{Latin Modern Math}
\usepackage{enumitem,needspace,xurl,hyperref}
\setlength{\parindent}{0pt}
\setlength{\parskip}{5pt}
\setlength{\emergencystretch}{3em}
\newcommand{\N}{\mathbb{N}}
\newcommand{\Z}{\mathbb{Z}}
\newcommand{\Q}{\mathbb{Q}}
\newcommand{\R}{\mathbb{R}}
\newcommand{\C}{\mathbb{C}}
\newcommand{\F}{\mathbb{F}}
\newcommand{\A}{\mathbb{A}}
\newcommand{\PP}{\mathbb{P}}
\newcommand{\E}{\mathbb{E}}
\newcommand{\eps}{\varepsilon}
\newcommand{\dd}{\,\mathrm{d}}
\begin{document}
\section*{UnsolvedMath ID: 30000759}
\textbf{Source problem code:} catalogue entry
\subsection{Definitions and mathematical statement}
Let $\mathfrak g$ be a finite-dimensional semisimple complex Lie algebra and let $\mathfrak p\subseteq\mathfrak g$ be a parabolic subalgebra containing a Borel subalgebra. Choose a Levi decomposition $\mathfrak p=\mathfrak a\ltimes\mathfrak u$, where $\mathfrak a$ is the reductive Levi part and $\mathfrak u$ the nilpotent radical. Let $V$ be a simple $\mathfrak a$-module. To form the indicated induced module, inflate $V$ to $\mathfrak p$ by making $\mathfrak u$ act trivially. With $U$ denoting universal enveloping algebra, put
\[
\Delta(\mathfrak p,V)=U(\mathfrak g)\otimes_{U(\mathfrak p)}V.
\]

For every such datum, the problem asks for the structure of $\Delta(\mathfrak p,V)$, in particular the isomorphism classes of its simple composition factors $L$ and their multiplicities $[\Delta(\mathfrak p,V):L]$. In a finite-length setting this means finding a composition series with simple quotients and determining the number of quotients of each isomorphism class, independent of the chosen series. These multiplicities describe composition data, not by themselves all extensions between factors.

The source allows an arbitrary simple Levi module: it does not assume that $V$ is finite dimensional, a highest-weight module, or in a specified block of category $\mathcal O$. Those restrictions are not added. The precise class in which a finite-length composition theory and an effective ``structure'' description are intended must therefore be clarified; no finite-length assertion is assumed solely from the extract. The supplied type-$A$ result concerns ``rough structure'', which is not defined here and is not identified with a complete fine composition or extension description.
\subsection{Short English statement}
Describe generalized Verma modules induced from arbitrary simple Levi modules, including all simple composition factors and multiplicities. The unrestricted inducing-module class and the distinction between the reported rough type-$A$ result and a full fine structure remain explicit.
\subsection{Sources}
\begin{itemize}
\item Supplied report: \emph{Algebraic Groups} (2008), Oberwolfach Reports, DOI 10.4171/owr/2007/22, pp. 1191--1242; proposed year 2007. \url{https://doi.org/10.4171/owr/2007/22}.
\item Supplied Mazorchuk--Stroppel reference on generalized Verma modules: \url{https://arxiv.org/abs/math/0702811}; supplied published \emph{Advances in Mathematics} version: \url{https://doi.org/10.1016/j.aim.2008.08.003}.
\item Local UnsolvedMath registry, permanent ID 30000759, catalogue entry; created 25 August 2026, updated 16 September 2026. Its corrected\_verified note says the full structure question was retained while notation was restored. The supplied 21 August 2026 literature triage labels the record partially solved, reporting a complete combinatorial description of rough structure in type $A$ while describing full fine composition structure for arbitrary generalized Verma modules as difficult and unresolved. This is a dated supplied assessment, not an independently checked current theorem statement.
\end{itemize}
\end{document}
